% GENERATED FILE. Edit canonical lessons, then run npm run generate:books.
% canonical-source-hash: fnv1a64:4507be74f7a91b8f
% canonical-lessons: RU-C284-intervyu, RU-C284-opros, RU-C284-redaktor, RU-C284-kommentirovat, RU-C284-tsitirovat

\chapter{Interviews and comments}
\label{ch:ru-interviews-and-comments}
\hlchaptermodality{284}

\begin{chapteropening}
\textbf{By the end of this chapter:} \emph{I can talk about an interview, a survey and an editor, and say that someone commented on or quoted something.}

Complete the last lesson of chapter 284: I can talk about an interview, a survey and an editor, and say that someone commented on or quoted something.
\end{chapteropening}

\section[\texorpdfstring{\ru{интервью} (intervyú)}{интервью (intervyú)}]{\ru{интервью} (intervyú) — an interview}

\label{lesson:RU-C284-intervyu}



\begin{quote}
Before the new one: say the Russian for to find out, then the Russian for to warn.
\end{quote}

\subsection*{You'll want to know: \ru{интервью}}
\textbf{\ru{интервью}} (\emph{intervyú}) — \textquotedblleft{}an interview\textquotedblright{}.

\begin{quote}
\ru{Певица} \ru{дала} \ru{интервью} \ru{газете}.
\end{quote}

\begin{quote}
\emph{Pevítsa dalá intervyú gazéte.}
\end{quote}

\begin{quote}
The singer gave an interview to a newspaper.
\end{quote}

Neuter and indeclinable: \emph{\ru{в} \ru{интервью}}. \emph{\ru{Брать} \ru{интервью} \ru{у}} + genitive — to interview someone.

\begin{cousinweb}
\emph{\ru{интервью}} was borrowed from English \emph{interview} (EtymDB), which comes from French \emph{entrevue} \textquotedblleft{}a meeting\textquotedblright{}, from the reflexive verb \emph{entrevoir} \textquotedblleft{}to see each other\textquotedblright{}.
\end{cousinweb}

\begin{grammarlens}[title={What you've built}]
One more word for telling, explaining and describing at B1.
\end{grammarlens}

\subsection*{Your turn}
\begin{itemize}
\raggedright
  \item \emph{Say it:} \emph{intervyú}
  \item \emph{Say it:} \emph{intervyú}, once more
  \item \emph{Say it:} a sentence of your own with \emph{intervyú}
  \item \emph{Recall:} say the Russian for to find out, then the Russian for to warn, then say \emph{intervyú} again
\end{itemize}

\begin{tcolorbox}[breakable,colback=teal!4,colframe=teal!35!black,title={Before you move on}]
What does \ru{интервью} mean? (\textquotedblleft{}An interview\textquotedblright{}.) Say it once more.
\end{tcolorbox}
\section[\texorpdfstring{\ru{опрос} (oprós)}{опрос (oprós)}]{\ru{опрос} (oprós) — a survey, a poll}

\label{lesson:RU-C284-opros}



\begin{quote}
Before the new one: say the Russian for to warn, then the Russian for an interview.
\end{quote}

\subsection*{You'll want to know: \ru{опрос}}
\textbf{\ru{опрос}} (\emph{oprós}) — \textquotedblleft{}a survey, a poll\textquotedblright{}.

\begin{quote}
\ru{После} \ru{опроса} \ru{все} \ru{узнали} \ru{мнение} \ru{студентов}.
\end{quote}

\begin{quote}
\emph{Pósle oprósa vse uználi mnéniye studéntov.}
\end{quote}

\begin{quote}
After the survey everyone learned what the students thought.
\end{quote}

\emph{\ru{Опрос} \ru{общественного} \ru{мнения}} — an opinion poll.

\begin{cousinweb}
\emph{\ru{опрос}} is built on \emph{\ru{о}-} and \emph{\ru{просить}} \textquotedblleft{}to ask\textquotedblright{}, which is inherited from Proto-Slavic \emph{prositi}, from PIE \emph{pre\'{k}-} \textquotedblleft{}to ask\textquotedblright{} (EtymDB): asking round.
\end{cousinweb}

\begin{grammarlens}[title={What you've built}]
One more word for telling, explaining and describing at B1.
\end{grammarlens}

\subsection*{Your turn}
\begin{itemize}
\raggedright
  \item \emph{Say it:} \emph{oprós}
  \item \emph{Say it:} \emph{oprós}, once more
  \item \emph{Say it:} a sentence of your own with \emph{oprós}
  \item \emph{Recall:} say the Russian for to warn, then the Russian for an interview, then say \emph{oprós} again
\end{itemize}

\begin{tcolorbox}[breakable,colback=teal!4,colframe=teal!35!black,title={Before you move on}]
What does \ru{опрос} mean? (\textquotedblleft{}A survey, a poll\textquotedblright{}.) Say it once more.
\end{tcolorbox}
\section[\texorpdfstring{\ru{редактор} (redáktor)}{редактор (redáktor)}]{\ru{редактор} (redáktor) — an editor}

\label{lesson:RU-C284-redaktor}



\begin{quote}
Before the new one: say the Russian for an interview, then the Russian for a survey.
\end{quote}

\subsection*{You'll want to know: \ru{редактор}}
\textbf{\ru{редактор}} (\emph{redáktor}) — \textquotedblleft{}an editor\textquotedblright{}.

\begin{quote}
\ru{Редактор} \ru{прочитал} \ru{статью} \ru{и} \ru{улыбнулся}.
\end{quote}

\begin{quote}
\emph{Redáktor prochitál statyú i ulybnúlsya.}
\end{quote}

\begin{quote}
The editor read the article and smiled.
\end{quote}

\emph{\ru{Главный} \ru{редактор}} — the editor-in-chief.

\begin{cousinweb}
EtymDB derives \emph{\ru{редактор}} from French \emph{rédacteur} \textquotedblleft{}editor\textquotedblright{}, given a Latin-style ending \emph{-\ru{ор}}; the French word is built on Latin \emph{redāctus}, the participle of \emph{redigere} \textquotedblleft{}to drive back, to bring back\textquotedblright{} (Lewis \& Short).
\end{cousinweb}

\begin{grammarlens}[title={What you've built}]
One more word for telling, explaining and describing at B1.
\end{grammarlens}

\subsection*{Your turn}
\begin{itemize}
\raggedright
  \item \emph{Say it:} \emph{redáktor}
  \item \emph{Say it:} \emph{redáktor}, once more
  \item \emph{Say it:} a sentence of your own with \emph{redáktor}
  \item \emph{Recall:} say the Russian for an interview, then the Russian for a survey, then say \emph{redáktor} again
\end{itemize}

\begin{tcolorbox}[breakable,colback=teal!4,colframe=teal!35!black,title={Before you move on}]
What does \ru{редактор} mean? (\textquotedblleft{}An editor\textquotedblright{}.) Say it once more.
\end{tcolorbox}
\section[\texorpdfstring{\ru{комментировать} (kommentírovat)}{комментировать (kommentírovat)}]{\ru{комментировать} (kommentírovat) — to comment on}

\label{lesson:RU-C284-kommentirovat}



\begin{quote}
Before the new one: say the Russian for a survey, then the Russian for an editor.
\end{quote}

\subsection*{You'll want to know: \ru{комментировать}}
\textbf{\ru{комментировать}} (\emph{kommentírovat}) — \textquotedblleft{}to comment on\textquotedblright{}.

\begin{quote}
\ru{Тренер} \ru{не} \ru{стал} \ru{комментировать} \ru{игру}.
\end{quote}

\begin{quote}
\emph{Tréner ne stal kommentírovat igrú.}
\end{quote}

\begin{quote}
The coach would not comment on the game.
\end{quote}

Imperfective; perfective \emph{\ru{прокомментировать}} (\emph{prokommentírovat}). It takes the accusative, with no preposition.

\begin{cousinweb}
\emph{\ru{комментировать}} is usually traced through German \emph{kommentieren} to Latin \emph{commentārī} \textquotedblleft{}to think over, to write about\textquotedblright{} (Lewis \& Short); EtymDB lists no parent for the Russian word.
\end{cousinweb}

\begin{grammarlens}[title={What you've built}]
One more word for telling, explaining and describing at B1.
\end{grammarlens}

\subsection*{Your turn}
\begin{itemize}
\raggedright
  \item \emph{Say it:} \emph{kommentírovat}
  \item \emph{Say it:} \emph{kommentírovat}, once more
  \item \emph{Say it:} a sentence of your own with \emph{kommentírovat}
  \item \emph{Recall:} say the Russian for a survey, then the Russian for an editor, then say \emph{kommentírovat} again
\end{itemize}

\begin{tcolorbox}[breakable,colback=teal!4,colframe=teal!35!black,title={Before you move on}]
What does \ru{комментировать} mean? (\textquotedblleft{}To comment on\textquotedblright{}.) Say it once more.
\end{tcolorbox}
\section[\texorpdfstring{\ru{цитировать} (tsitírovat)}{цитировать (tsitírovat)}]{\ru{цитировать} (tsitírovat) — to quote}

\label{lesson:RU-C284-tsitirovat}



\begin{quote}
Before the new one: say the Russian for an editor, then the Russian for to comment on.
\end{quote}

\subsection*{You'll want to know: \ru{цитировать}}
\textbf{\ru{цитировать}} (\emph{tsitírovat}) — \textquotedblleft{}to quote\textquotedblright{}.

\begin{quote}
\ru{Учитель} \ru{любил} \ru{цитировать} \ru{стихи}.
\end{quote}

\begin{quote}
\emph{Uchítel lyubíl tsitírovat stikhí.}
\end{quote}

\begin{quote}
The teacher liked to quote poems.
\end{quote}

Imperfective; perfective \emph{\ru{процитировать}} (\emph{protsitírovat}). \emph{\ru{Цитата}} — a quotation.

\begin{cousinweb}
\emph{\ru{цитировать}} was borrowed from German \emph{zitieren} with the Russian ending \emph{-\ru{овать}} (EtymDB); the German verb comes from Latin \emph{citāre} \textquotedblleft{}to set in motion, to summon\textquotedblright{} (Lewis \& Short).
\end{cousinweb}

\begin{grammarlens}[title={What you've built}]
One more word for telling, explaining and describing at B1.
\end{grammarlens}

\subsection*{Your turn}
\begin{itemize}
\raggedright
  \item \emph{Say it:} \emph{tsitírovat}
  \item \emph{Say it:} \emph{tsitírovat}, once more
  \item \emph{Say it:} a sentence of your own with \emph{tsitírovat}
  \item \emph{Recall:} say the Russian for an editor, then the Russian for to comment on, then say \emph{tsitírovat} again
\end{itemize}

\begin{tcolorbox}[breakable,colback=teal!4,colframe=teal!35!black,title={Before you move on}]
What does \ru{цитировать} mean? (\textquotedblleft{}To quote\textquotedblright{}.) Say it once more.
\end{tcolorbox}
